% Source ownership:
% - Manel: collaborative equilibration, Boltzmann roulette, metrics
% - Arthur: ensemble averaging / orchestrator narrative
% - Itxaso: MUST add "Orchestrator Initial Config Read & Broadcast"

\section{MPI Parallelization: Collaborative Equilibration}
\label{sec:mpi_collab}

\subsection{Orchestrator Initial Configuration Read and Broadcast}

In the MPI version of the program, rank 0 acts as the central
orchestrator. Its role is to read the initial configuration files,
decide whether each conformation must still be equilibrated or can
directly enter production, and distribute the corresponding work to the
workers.

At startup, the master reads the input settings and the geometry
required to initialise each polymer configuration. If an equilibrated
geometry is already available for a given configuration, the program can
bypass the equilibration stage and directly unlock the corresponding
production replicas. Otherwise, the master marks that configuration as
pending collaborative equilibration.

For configurations that still require equilibration, rank 0 assigns a
group of workers to the same initial structure. These workers all begin
from the same geometry, but each one uses a different random seed so
their trajectories immediately diverge. The master therefore provides a
common starting point while still allowing independent exploration of
phase space.

This orchestrator stage also centralises the campaign-level file
handling. Workers do not need to parse the full simulation manifest or
determine which configurations are ready for production. Instead, they
receive either the job descriptor or the coordinates required for the
task they have been assigned, initialise their local arrays, and begin
Monte Carlo propagation.

\subsection{Equilibration Optimization}

Rather than assigning a single worker to equilibrate each configuration
type independently, a group of \texttt{workers\_per\_equil} MPI workers
is assigned to each configuration type simultaneously. Each worker
starts from the same initial geometry but with a different random seed,
for example \texttt{rngseed = rank * 104729}, ensuring that trajectories
immediately diverge and explore distinct regions of conformational
space.

Every \texttt{sync\_interval} Monte Carlo steps, a synchronisation
protocol is executed. Each worker sends its current energy and
\(R_g\) to the master, which selects the most representative
configuration via a Boltzmann-weighted roulette wheel and broadcasts
those coordinates back to the non-winning workers in the group. The
winning worker continues its trajectory uninterrupted, preserving the
best conformational path found so far.

The purpose of this strategy is to reduce the time required to reach a
well-equilibrated low-energy configuration. Instead of waiting for one
trajectory to escape a poor metastable region on its own, the algorithm
lets several replicas search in parallel and periodically shares the
best result across the group. In this way, all workers benefit from the
progress made by the most successful replica.

\subsection{Replica Selection: Boltzmann Roulette Wheel}

At each synchronisation point, the master assigns a statistical weight
to each replica \(i\) based on its total energy \(E_i\):
\[
w_i = \exp\!\left(-\frac{E_i - E_{\min}}{k_B T_{\mathrm{virt}}}\right),
\]
where \(E_{\min}\) is the lowest energy in the current group and
\(T_{\mathrm{virt}}\) is a virtual temperature.

The virtual temperature is a purely algorithmic parameter with no
physical meaning. It controls how aggressively low-energy replicas are
favoured:
\begin{itemize}
\item \(T_{\mathrm{virt}} \to 0\): deterministic selection of the
minimum-energy replica.
\item \(T_{\mathrm{virt}} \to \infty\): uniform selection, ignoring
energies.
\end{itemize}

A uniform random number then samples this distribution via roulette-wheel
selection. The selected replica becomes the winner of that
synchronisation round, and its coordinates are sent to the non-winning
workers so they can resume the search from the most promising state
found so far.

\subsection{MPI Communication Protocol}

The synchronisation uses only point-to-point MPI operations
\texttt{MPI\_Send} and \texttt{MPI\_Recv} to avoid the need for
temporary communicators or collective operations across worker
subgroups. At each sync point, the master loops over all workers in the
active equilibration group and performs the following exchanges.

\begin{table}[h!]
\centering
\begin{tabular}{clll}
\toprule
Step & Direction & Tag & Content \\
\midrule
1 & worker \(\rightarrow\) master & \texttt{TAG\_SYNC\_OBS} &
\(E_{\mathrm{tot}}, R_g^2\) \\
2 & master \(\rightarrow\) worker & \texttt{TAG\_SYNC\_CTRL} &
control signal, \texttt{bestlocalidx} \\
3 & worker \(\rightarrow\) master & \texttt{TAG\_EQUIL\_COORDS} &
full coordinate array \\
4 & master \(\rightarrow\) non-winners & \texttt{TAG\_SYNC\_COORDS} &
winning coordinates \\
\bottomrule
\end{tabular}
\caption{Point-to-point synchronisation protocol used during collaborative equilibration.}
\end{table}

Once the Geweke test declares equilibrium, the master saves the
equilibrated geometry to \texttt{../results/equilibrated\_cX.xyz} and
unlocks the 10 production jobs for that configuration type in the task
queue.

\subsection{Metrics and Methodology}

This section evaluates the performance of the equilibration
optimization technique compared to the serial baseline. The system
studied is configuration 1, a chain of 500 carbons at \(T = 300\) K
including hydrogens. Efficiency is evaluated based on the time to reach
specific energy thresholds \(E_{\mathrm{tot}} = 110\), 100, and
90 kcal/mol.

The speedup is defined as
\[
S_w = \frac{\mathrm{Steps}_{\mathrm{serial}}}
           {\mathrm{Steps}_{\mathrm{parallel},\,w}},
\]
where \(w\) is the number of workers participating in the collaborative
equilibration group. Efficiency is then defined as
\[
\eta_w = \frac{S_w}{w}.
\]

The table below summarises the number of MC steps required by different
worker configurations to reach the target energy levels.

\begin{table}[h!]
\centering
\begin{tabular}{ccccc}
\toprule
Threshold & Workers \(w\) & Steps Required & Speedup \(S_w\) & Efficiency \(S_w / w\) \\
\midrule
110 kcal/mol & 1 & 590,000  & 1.00  & 1.00 \\
110 kcal/mol & 2 & 70,000   & 8.43  & 4.21 \\
110 kcal/mol & 4 & 140,000  & 4.21  & 1.05 \\
110 kcal/mol & 8 & 50,000   & 11.80 & 1.47 \\
\midrule
100 kcal/mol & 1 & 1,180,000 & 1.00  & 1.00 \\
100 kcal/mol & 2 & 350,000   & 3.37  & 1.68 \\
100 kcal/mol & 4 & 210,000   & 5.62  & 1.40 \\
100 kcal/mol & 8 & 50,000    & 23.60 & 2.95 \\
\midrule
90 kcal/mol & 1 & 2,400,000 & 1.00  & 1.00 \\
90 kcal/mol & 2 & 1,780,000 & 1.35  & 0.67 \\
90 kcal/mol & 4 & 570,000   & 4.21  & 1.05 \\
90 kcal/mol & 8 & 160,000   & 15.00 & 1.87 \\
\bottomrule
\end{tabular}
\caption{Monte Carlo steps required to reach selected energy thresholds during collaborative equilibration.}
\end{table}

These results show that the method can exhibit superlinear speedup.
This is not just because more workers are used, but because the search
itself changes. A population of replicas explores several
conformational regions at once, and once one replica finds a much
better basin, that discovery is broadcast to the others at the next
synchronisation point.

For this reason, the observed acceleration is partly computational and
partly algorithmic. The collaborative protocol does not simply divide
the same work among more processes; it improves the probability of
finding favourable states earlier than a single serial trajectory would.

\begin{figure}[h!]
  \centering
  \includegraphics[width=\textwidth]{convergence_speedup_equil.pdf}
  \caption{Convergence speedup analysis for collaborative equilibration.}
\end{figure}

\subsection{Ensemble Averaging and Star Topology}

Once an equilibrated geometry has been found, the workflow switches from
collaborative equilibration to large-scale production sampling. At this
stage, rank 0 acts as a central dispatcher in a star-topology MPI
scheme, while the remaining ranks behave as workers that repeatedly
request new tasks from the master.

The master maintains a global queue of production jobs. Each job
corresponds to an independent production replica, and workers are
assigned new jobs dynamically as soon as they finish their current one.
This avoids the load imbalance that would appear if all replicas were
distributed statically at launch, since individual trajectories may
still differ slightly in runtime.

The dispatcher listens using a tag-based communication pattern and can
react to whichever worker becomes free first. In practice, this means
that workers request work, receive a production assignment, run the
simulation, write their outputs, and notify the master when they have
finished. The master then either sends another job or terminates the
worker once the queue has been exhausted.

This star-topology design turns the MPI layer into an ensemble averaging
engine. Independent production replicas are generated as efficiently as
possible, and the final observables are obtained by averaging over the
outputs of all completed runs. In this phase, the objective is no
longer to accelerate the equilibration of a single chain, but to
maximize statistical throughput across the whole ensemble.

\begin{table}[h!]
\centering
\begin{tabular}{cll}
\toprule
Direction & Tag & Meaning \\
\midrule
worker \(\rightarrow\) master & \texttt{TAG\_REQUEST\_WORK} & worker requests a new task \\
master \(\rightarrow\) worker & \texttt{TAG\_DO\_PROD} & start a production replica \\
master \(\rightarrow\) worker & \texttt{TAG\_WAIT} & remain idle temporarily \\
master \(\rightarrow\) worker & \texttt{TAG\_DIE} & terminate execution \\
worker \(\rightarrow\) master & \texttt{TAG\_PROD\_DONE} & production replica finished \\
\bottomrule
\end{tabular}
\caption{Tag-based work distribution in the star-topology production phase.}
\end{table}

\begin{figure}[h!]
  \centering
  \includegraphics[width=\textwidth]{parallel_star_speedup_plot.png}
  \caption{Timing and efficiency of the orchestrator-worker production workflow.}
\end{figure}